% =====================================================================
% Mirante dos Dados — Working Paper #4
% Análise Estatística da Prevalência da Doença de Parkinson no Brasil:
% Desafios, Tecnologias de Neuroimagem e Perspectivas de Políticas Públicas
%
% Autores: Alexandre Maciel Rolim (manuscrito original — epidemiologia,
%          revisão de literatura, recomendações de política pública)
%          + Leonardo Chalhoub (camada de engenharia de dados, análise
%          atualizada de infraestrutura via microdados CNES)
% Compilação: pdflatex articles/equipamentos-rm-parkinson.tex (2x)
% =====================================================================

\documentclass[12pt,a4paper,oneside]{article}

\usepackage[T1]{fontenc}
\usepackage[utf8]{inputenc}
\usepackage[brazilian]{babel}
\usepackage{lmodern}
\usepackage[a4paper,top=3cm,bottom=2cm,left=3cm,right=2cm]{geometry}
\usepackage{setspace}
\usepackage{indentfirst}
\usepackage{titlesec}
\usepackage{caption}
\usepackage{booktabs}
\usepackage{array}
\usepackage{multirow}
\usepackage{graphicx}
\usepackage[hidelinks]{hyperref}
\usepackage{url}
\usepackage{float}
\usepackage{tcolorbox}

\graphicspath{{figures-equipamentos-rm/}}

\onehalfspacing
\setlength{\parindent}{1.25cm}
\setlength{\parskip}{0pt}

\titleformat{\section}{\normalfont\bfseries\large\uppercase}{\thesection}{0.6em}{}
\titleformat{\subsection}{\normalfont\bfseries\normalsize}{\thesubsection}{0.6em}{}
\titlespacing*{\section}{0pt}{18pt}{8pt}
\titlespacing*{\subsection}{0pt}{12pt}{6pt}

\captionsetup[table]{labelfont=bf,name=Tabela,labelsep=endash,
                     position=above,justification=centering,font=small}
\captionsetup[figure]{labelfont=bf,name=Figura,labelsep=endash,
                      position=below,justification=centering,font=small}

\newcommand{\rs}{R\$\,}

\begin{document}

% ---------- CAPA ---------------------------------------------------------
\begin{titlepage}
  \begin{center}
    \vspace*{1.5cm}
    {\small\textsc{Mirante dos Dados \quad\textbar\quad Working Paper n. 4}}\\
    {\small Versão 2.0 \textemdash{} Abril de 2026}

    \vfill

    {\LARGE\bfseries
     ANÁLISE ESTATÍSTICA DA PREVALÊNCIA DA DOENÇA DE
     PARKINSON NO BRASIL: DESAFIOS, TECNOLOGIAS DE
     NEUROIMAGEM E PERSPECTIVAS DE POLÍTICAS PÚBLICAS}

    \vspace{1cm}
    {\large\itshape
     Statistical analysis of Parkinson's disease prevalence in Brazil:
     challenges, neuroimaging technologies, and public-policy perspectives}

    \vfill

    {\large\textbf{Alexandre Maciel Rolim}\textsuperscript{1} \quad\textbullet\quad
     \textbf{Leonardo Chalhoub}\textsuperscript{2}}\\[6pt]
    {\small\textsuperscript{1}\,Pesquisa original \textemdash{} epidemiologia, revisão
     de literatura, síntese de políticas públicas (manuscrito v1, abr/2026).}\\[2pt]
    {\small\textsuperscript{2}\,Engenharia de dados \textemdash{} pipeline aberto e
     análise atualizada de infraestrutura via microdados CNES 2013\textendash 2025.}\\[6pt]
    {\small\href{https://github.com/leonardochalhoub/mirante-dos-dados-br}{github.com/leonardochalhoub/mirante-dos-dados-br}}

    \vfill

    {\small Brasil \\ Abril de 2026}
  \end{center}
\end{titlepage}

% ---------- RESUMO -------------------------------------------------------
\section*{Resumo}
\addcontentsline{toc}{section}{Resumo}
\begin{singlespace}\small\noindent
A doença de Parkinson (DP) é a segunda condição neurodegenerativa mais
prevalente globalmente, afetando aproximadamente 8,5 milhões de
indivíduos com projeções de 25,2 milhões de casos até 2050. Este artigo
integra a análise epidemiológica do manuscrito original de Rolim (2026)
\textemdash{} fundamentado nos achados do estudo ELSI-Brazil
(2019\textendash 2021) que estimou prevalência de 0,84\,\%
(IC 95\,\%: 0,64\textendash 1,09) em brasileiros com 50 anos ou mais
e projetou 1.250.638 casos de DP no Brasil até 2060 \textemdash{} com
uma análise atualizada e reproduzível da infraestrutura de neuroimagem
no SUS, derivada dos microdados CNES (DATASUS) processados pelo
pipeline aberto do projeto \textit{Mirante dos Dados}. Documentamos
seis achados principais. \textbf{Primeiro}, o Brasil dispunha em 2025
de aproximadamente 10.079 aparelhos de Ressonância Magnética
distribuídos em 6.744 estabelecimentos, com participação do SUS
estável em torno de 43\,\% (4.317 aparelhos). \textbf{Segundo}, a
densidade nacional alcançou 47,2 aparelhos por milhão de habitantes,
patamar superior à mediana da OCDE (\,$\sim$\,17/Mhab em 2021), mas
essa média mascara desigualdade regional severa. \textbf{Terceiro}, o
gradiente regional é nítido: Sudeste 64,82/Mhab vs.\ Norte 25,55/Mhab,
razão de 2,54:1, com Sudeste concentrando 57\,\% do parque nacional
para 42\,\% da população. \textbf{Quarto}, o coeficiente de variação
per capita entre as 27 UFs permaneceu estável em 0,44\textendash 0,49
ao longo de 2013\textendash 2025, indicando que a expansão acelerada
do parque (de 22,5 para 47,2/Mhab, +110\,\%) NÃO reduziu a iniquidade
geográfica relativa. \textbf{Quinto}, o estado do Rio de Janeiro
lidera em densidade (82,99/Mhab), enquanto Amazonas, Roraima, Piauí
e Acre ficam abaixo da metade da mediana OCDE (15,5\textendash 22,6/Mhab),
configurando situação de subdotação clinicamente preocupante para o
diagnóstico diferencial de DP em estágios iniciais. \textbf{Sexto}, a
implementação seletiva de protocolos padronizados de RM em centros
de referência regional, conforme proposto por Rolim (2026), exigiria
investimento único estimado em R\$\,500 mil\textendash 1 milhão e
operaria dentro das restrições orçamentárias realistas do SUS. O
trabalho contribui ao ecossistema de pesquisa em saúde pública
brasileira ao disponibilizar (a) o manuscrito clínico-epidemiológico
de Rolim como referência base e (b) a infraestrutura aberta do
Mirante dos Dados para reanálise contínua e integração com novos
recortes (e.g., Bolsa Família, Emendas Parlamentares) sem custo
marginal de processamento de microdados.

\vspace{0.4cm}\noindent\textbf{Palavras-chave:}
doença de Parkinson; ELSI-Brazil; ressonância magnética; SUS; CNES;
neuroimagem; saúde pública; iniquidade regional; dados abertos.
\end{singlespace}

% ---------- ABSTRACT -----------------------------------------------------
\section*{Abstract}
\addcontentsline{toc}{section}{Abstract}
\begin{singlespace}\small\noindent
Parkinson's disease (PD) is the second most prevalent
neurodegenerative condition globally, affecting approximately
8.5 million individuals, with projections of 25.2 million cases by
2050. This paper integrates the epidemiological analysis from the
original manuscript by Rolim (2026) \textemdash{} grounded in the
ELSI-Brazil study (2019\textendash 2021) which estimated PD
prevalence of 0.84\% (95\% CI 0.64\textendash 1.09) among Brazilians
aged 50+ and projected 1,250,638 PD cases in Brazil by 2060
\textemdash{} with an updated, reproducible analysis of MRI
infrastructure across the Brazilian Unified Health System (SUS),
derived from CNES (DATASUS) microdata processed by the open
\textit{Mirante dos Dados} pipeline. We document six main findings.
\textbf{First}, Brazil had \,$\sim$\,10,079 MRI scanners across 6,744
facilities in 2025, with SUS share stable around 43\%.
\textbf{Second}, national density reached 47.2/Mpop (above OECD
median \,$\sim$\,17), but this average masks severe regional
inequality. \textbf{Third}, the regional gradient is sharp: Southeast
64.82/Mpop vs.\ North 25.55/Mpop. \textbf{Fourth}, per-capita
coefficient of variation across UFs remained stable at 0.44\textendash 0.49
during 2013\textendash 2025 \textemdash{} the rapid expansion did NOT
reduce relative geographic inequality. \textbf{Fifth}, Amazonas,
Roraima, Piauí and Acre fall below half the OECD median, posing a
clinical concern for early-stage PD differential diagnosis.
\textbf{Sixth}, Rolim's selective implementation of standardized MRI
protocols at regional reference centers would cost an estimated
BRL 500K\textendash 1M one-time, fitting realistic SUS budgetary
constraints.

\vspace{0.4cm}\noindent\textbf{Keywords:}
Parkinson's disease; ELSI-Brazil; magnetic resonance imaging; SUS;
CNES; neuroimaging; public health; regional inequality; open data.
\end{singlespace}

\newpage

% ---------- SUMÁRIO + LISTAS --------------------------------------------
\renewcommand{\contentsname}{Sumário}
\tableofcontents

\newpage
\renewcommand{\listfigurename}{Lista de Figuras}
\listoffigures

\vspace{1cm}
\renewcommand{\listtablename}{Lista de Tabelas}
\listoftables

\newpage

% =====================================================================
\section{Introdução}\label{sec:intro}
% =====================================================================

A doença de Parkinson constitui a segunda condição neurodegenerativa
mais prevalente globalmente, afetando aproximadamente 8,5 milhões de
indivíduos\,[\ref{ref:elsi}], com projeções estimando aumento para
25,2 milhões de casos até 2050\,[\ref{ref:su}]. No Brasil, um país de
renda média-alta com população em transição demográfica acelerada, a
DP representa desafio crescente para sistemas de saúde pública,
especialmente considerando a heterogeneidade regional e as
disparidades em acesso a recursos diagnósticos especializados.
Estudos epidemiológicos recentes, particularmente o ELSI-Brazil
(Estudo Longitudinal da Saúde dos Idosos Brasileiros), revelaram que
a prevalência de DP entre indivíduos com 50 anos ou mais é de 0,84\%
(IC 95\,\%: 0,64\,\%\textendash 1,09\,\%), com prevalência
padronizada por idade e sexo de 0,85\,\% (IC 95\,\%:
0,62\,\%\textendash 1,10\,\%)\,[\ref{ref:elsi}].

A relevância de estudar a prevalência de DP no Brasil transcende
meramente a quantificação de casos: implica compreender os
mecanismos pelos quais o diagnóstico é realizado, os desafios
tecnológicos enfrentados pelo sistema público de saúde, e as
oportunidades de implementação de políticas públicas que integrem
epidemiologia com tecnologias diagnósticas avançadas. A ressonância
magnética, enquanto ferramenta não-invasiva de neuroimagem
estrutural e funcional, emerge como instrumento potencial para
melhorar acurácia diagnóstica em contexto brasileiro, permitindo
identificação de características anatômicas e funcionais associadas
a Parkinson e condições que mimetizam a doença\,[\ref{ref:saba}].

\subsection{Origem e estrutura deste artigo}

Este Working Paper integra duas contribuições. A primeira, do autor
\textbf{Alexandre Maciel Rolim}, é o manuscrito clínico-epidemiológico
original (v1, abril de 2026) \textemdash{} integrando análise
epidemiológica do ELSI-Brazil com revisão de tecnologias de
neuroimagem e síntese de recomendações de política pública. A segunda, do autor
\textbf{Leonardo Chalhoub}, é a camada de engenharia de dados:
pipeline aberto do projeto \textit{Mirante dos Dados} que processa
microdados mensais do CNES (DATASUS) para 2005\textendash 2025,
gerando indicadores reproduzíveis de infraestrutura de Ressonância
Magnética por UF, ano e setor (SUS vs.\ Privado). A primeira
contribuição fornece o arcabouço epidemiológico e o aparato
clínico\textbf{; a segunda atualiza, formaliza e torna reproduzível
o panorama da infraestrutura} \textemdash{} oferecendo, para cada
afirmação de caráter quantitativo, dados auditáveis e cobertura
temporal estendida.

A linha do tempo do contexto clínico-tecnológico analisado é
sintetizada na Figura\,\ref{fig:timeline}.

\begin{figure}[H]
\centering
\includegraphics[width=0.95\linewidth]{fig01-timeline-rm}
\caption{Linha do tempo do contexto da pesquisa em DP e RM no Brasil}
\label{fig:timeline}
\par\noindent{\small Fonte: elaboração própria. CNES inicia em 2005;
SIGTAP padronizado em 2008; descrição do \textit{swallow-tail sign}
(Schwarz \textit{et al.}, 2014); critérios diagnósticos clínicos
MDS-PD (Postuma \textit{et al.}, 2015); imagem com sensibilidade a
neuromelanina (Pyatigorskaya \textit{et al.}, 2018).}
\end{figure}

% =====================================================================
\section{Revisão de Literatura}\label{sec:teoria}
% =====================================================================

\subsection{Epidemiologia da Doença de Parkinson no Brasil:
achados recentes e tendências}

A epidemiologia de Parkinson no Brasil permaneceu incompletamente
caracterizada até recentemente. O único estudo populacional anterior
com representatividade metodológica foi o estudo Bambuí, publicado
em 2006, que identificou prevalência de 3,3\% em população acima de
65 anos em município único\,[\ref{ref:elsi}]. A lacuna de dezoito
anos sem dados epidemiológicos nacionais abrangentes constituía
obstáculo crítico para planejamento adequado de políticas de saúde.
O ELSI-Brazil, coletando dados entre 2019 e 2021 de 9.949
participantes distribuídos por 70 municípios nas cinco regiões
geopolíticas brasileiras, forneceu primeira estimativa
representativa nacionalmente. A prevalência identificada de 0,84\%
entre brasileiros com 50 anos ou mais alinha-se com estimativas de
outros países de renda média, como China (0,94\%)\,[\ref{ref:elsi}],
embora significativamente inferior ao relatado no estudo Bambuí,
sugerindo possível subestimação contemporânea relacionada a
diagnóstico autorrelatado, redução verdadeira em prevalência
relativa, ou diferenças metodológicas entre estudos.

A distribuição etária da prevalência em ELSI-Brazil revelou padrão
exponencial esperado: 0,39\% (IC 95\,\%: 0,17\%\textendash 0,60\%)
nos 50\textendash 59 anos, elevando-se para 2,75\% (IC 95\,\%:
1,39\%\textendash 4,11\%) em indivíduos com 80 anos ou
mais\,[\ref{ref:elsi}]. Esta progressão reflete fenômeno
bem-documentado de que envelhecimento é fator de risco primário
para Parkinson, relacionado a acúmulo progressivo de patologia de
alfa-sinucleína e morte neuronal dopaminérgica ao longo do
tempo\,[\ref{ref:bovolenta}]. A prevalência foi significativamente
superior em homens (1,21\%, IC 95\,\%: 0,79\%\textendash 1,63\%)
comparada a mulheres (0,52\%, IC 95\,\%: 0,31\%\textendash 0,73\%),
com razão de chances de 2,35 (IC 95\,\%:
1,35\textendash 4,08; p\,=\,0,003)\,[\ref{ref:elsi}], alinhando-se
com literatura global documentando dimorfismo sexual pronunciado
em Parkinson.

\subsection{Neuroimagem estrutural e funcional na detecção de DP}

A ressonância magnética estrutural constitui ferramenta diagnóstica
potencialmente valiosa em Parkinson, permitindo avaliação de
integridade neuronal, volume de substância negra, e padrões de
atrofia regional\,[\ref{ref:saba}]. Estudos demonstram que
volumetria da substância negra, quantificada por técnicas de
segmentação automatizada em imagens de RM, correlaciona com grau
de deficiência dopaminérgica avaliada por single-photon emission
computed tomography (SPECT) com ligante de transportador de
dopamina (DAT-SPECT)\,[\ref{ref:saba}]. A espectroscopia de
ressonância magnética (MRS), técnica que mede concentrações de
metabolitos cerebrais, permite quantificação de glutamato, GABA,
lactato, e outros metabólitos que frequentemente mostram alterações
em Parkinson\,[\ref{ref:saba}].

Estudos de RM funcional de repouso demonstram alterações em
conectividade funcional de redes cerebrais em indivíduos com
Parkinson comparados a controles, particularmente em redes de modo
padrão (default mode network)\,[\ref{ref:saba}]. Ressonância
magnética de alta resolução com sequências T1 e T2 permite
identificação de hiper-intensidades em núcleo rubro, putame, e
outras estruturas de gânglios basais que auxiliam em diferenciação
entre Parkinson idiopático e parkinsonismos atípicos (atrofia
multissistêmica, paralisia supranuclear progressiva)\,[\ref{ref:saba}].

\subsection{Desafios de infraestrutura de neuroimagem no SUS}

O Sistema Único de Saúde brasileiro enfrenta desafios substanciais
na provisão equitativa de serviços de neuroimagem diagnóstica. A
distribuição de equipamentos de RM no Brasil é profundamente
desigual: a região Sudeste concentra parcela majoritária do parque
nacional, enquanto a região Norte permanece severamente
subdotada\,[\ref{ref:datasus}]. Esta distribuição geográfica reflete
padrões históricos de desenvolvimento econômico desigual que
criaram disparidades estruturais em acesso a tecnologia
diagnóstica. Para DP especificamente, a implicação é que indivíduos
em regiões periféricas com apresentação clínica ambígua enfrentam
barreiras praticamente intransponíveis para acesso a neuroimagem
especializadas que poderiam confirmar ou excluir diagnóstico.

Adicionalmente, o custo de manutenção, calibração, e atualização de
equipamentos de RM é extraordinário, com custos operacionais anuais
de máquina clínica de 3 Tesla variando de \rs 200.000 a
\rs 400.000 por ano\,[\ref{ref:datasus}]. Para países de renda média
como Brasil, estes custos representam proporção significativa de
orçamentos de saúde, criando trade-offs entre investimento em
neuroimagem avançada versus alocação de recursos para outras
prioridades de saúde pública.

% =====================================================================
\section{Aspectos Metodológicos}\label{sec:metodo}
% =====================================================================

\subsection{Desenho de estudo e fontes de dados}

Este artigo constitui revisão sistemática qualitativa de dados
epidemiológicos de prevalência de DP no Brasil
\textemdash{} \textbf{contribuição original de Rolim (2026)}
\textemdash{} integrada com análise quantitativa atualizada e
reproduzível de infraestrutura tecnológica disponível no SUS
\textemdash{} \textbf{contribuição original de Chalhoub (2026)}
via pipeline \textit{Mirante dos Dados}.

Não foi realizada experimentação com seres humanos ou animais; a
base de dados utilizada compreende fontes secundárias publicamente
disponíveis: \textbf{(1)} estudo ELSI-Brazil, segunda onda
(2019\textendash 2021), que forneceu dados epidemiológicos de
prevalência, distribuição etária-sexual, e características clínicas
associadas de DP em população representativa
brasileira\,[\ref{ref:elsi}]; \textbf{(2)} microdados mensais do
\textbf{Cadastro Nacional de Estabelecimentos de Saúde (CNES)},
mantidos pelo DATASUS, contendo cadastro de cada equipamento de RM
por estabelecimento, UF, mês, ano e indicador SUS/Privado; foram
processados aproximadamente 6.614 arquivos \texttt{.dbc}
correspondentes a 27 UFs $\times$ 246 meses
(jan/2005\textendash dez/2025); \textbf{(3)} estimativas anuais de
população residente por UF, da Tabela 6579 do IBGE/SIDRA, para
normalização per capita; \textbf{(4)} literatura crítica sobre
tecnologias de neuroimagem em DP (PubMed, Scopus, Web of Science,
2015\textendash 2024).

\subsection{Pipeline de dados}

A camada de engenharia de dados adota a arquitetura medallion
(ARMBRUST \textit{et al.}, 2021) implementada em Apache Spark e
persistida em Delta Lake. A síntese arquitetural é apresentada na
Figura\,\ref{fig:architecture}.

\begin{figure}[H]
\centering
\includegraphics[width=0.95\linewidth]{fig02-architecture}
\caption{Pipeline medallion (bronze\,/\,silver\,/\,gold) usado para
materializar o painel CNES de Ressonância Magnética}
\label{fig:architecture}
\par\noindent{\small Fonte: elaboração própria, baseada em
Armbrust \textit{et al.} (2021).}
\end{figure}

\textbf{Bronze} converte cada arquivo DBC para Parquet (via PySUS) e
materializa Delta com partição UF$\times$ano. \textbf{Silver} agrega
por (UF, ano, codequip), preservando todos os 99 códigos de
equipamento mas separando contagens SUS/Privado. \textbf{Gold}
colapsa para o painel UF$\times$ano$\times$codequip e adiciona
população IBGE para normalização per capita. Para o presente artigo,
o filtro \texttt{codequip = 42} extrai apenas Ressonância Magnética.
O pipeline está hospedado em Databricks Free Edition com refresh
mensal automatizado.

\subsection{Indicadores construídos}

Para cada (UF, ano), o painel disponibiliza: \textbf{(i)} contagem
total média mensal de aparelhos de RM; \textbf{(ii)} divisão
SUS/Privado; \textbf{(iii)} número de estabelecimentos distintos com
$\geq$\,1 RM; \textbf{(iv)} densidade per capita (RM/Mhab) usando
população IBGE-SIDRA; \textbf{(v)} série temporal completa
2013\textendash 2025 (a janela 2005\textendash 2012 está
materializada na bronze mas não foi usada neste artigo por
limitações de dimensionamento da análise comparativa).

\subsection{Limitações}

\textbf{Primeira}, a contagem CNES reflete equipamentos
\textit{cadastrados}, não \textit{operacionais} \textemdash{} a
diferença entre cadastro e operação efetiva é não-trivial e exigiria
cruzamento com produção (SIH-RD/SIA-RD) para mensuração rigorosa.
\textbf{Segunda}, o cadastro pode duplicar equipamentos quando o
mesmo aparelho é registrado por múltiplos profissionais
prestadores; controles foram aplicados via \texttt{COUNT(DISTINCT
CNES)} mas resíduos podem persistir. \textbf{Terceira}, a análise é
descritiva: não pretende estabelecer relação causal entre
disponibilidade de RM e desfechos clínicos de DP. \textbf{Quarta}, o
manuscrito original de Rolim (2026) usa estimativa de
\,$\sim$\,1.800 equipamentos no país (referência DATASUS-2024),
enquanto nossa contagem para 2025 é de 10.079; a diferença reflete
metodologia distinta (cadastro tipo único vs.\ contagem de unidades
por estabelecimento) e ressalta a importância de se trabalhar com
microdados auditáveis.

% =====================================================================
\section{Resultados}\label{sec:resultados}
% =====================================================================

\subsection{Volumes nacionais e evolução temporal}

A Figura\,\ref{fig:evolution-total} apresenta a evolução do parque
total de aparelhos de RM no Brasil entre 2013 e 2025. O volume mais
que dobrou no período: de 4.526 para 10.079 unidades (+123\%), com
crescimento per capita de 22,5 para 47,2 aparelhos por milhão de
habitantes (+110\%, refletindo a expansão líquida descontada do
crescimento populacional).

\begin{figure}[H]
\centering
\includegraphics[width=0.95\linewidth]{fig03-evolution-total}
\caption{Brasil — número total de aparelhos de RM (média mensal por ano), 2013\textendash 2025}
\label{fig:evolution-total}
\par\noindent{\small Fonte: elaboração própria, microdados CNES/DATASUS.}
\end{figure}

A Figura\,\ref{fig:evolution-pm} expressa a mesma trajetória em
densidade por milhão de habitantes, com linha de referência da
mediana OCDE ($\sim$\,17/Mhab, OECD Health Statistics
2021\textendash 2023). \textbf{Brasil ultrapassou a mediana OCDE em
todos os anos da série}; em 2025 chegou a 47,2/Mhab, próximo da
média de países como Coreia do Sul (\,$\sim$\,30/Mhab) e Itália
(\,$\sim$\,30/Mhab), mas ainda abaixo de Japão (\,$\sim$\,57/Mhab) e
EUA (\,$\sim$\,38/Mhab).

\begin{figure}[H]
\centering
\includegraphics[width=0.95\linewidth]{fig04-evolution-per-million}
\caption{Densidade nacional de RM por milhão de habitantes,
2013\textendash 2025, com referência da mediana OCDE}
\label{fig:evolution-pm}
\par\noindent{\small Fonte: elaboração própria, CNES/DATASUS + IBGE/SIDRA.}
\end{figure}

A Tabela\,\ref{tab:evolucao-rm} consolida a série temporal completa.

\begin{table}[H]
\centering
\caption{Aparelhos de RM no Brasil por ano e setor, 2013\textendash 2025}
\begin{tabular}{rrrrrr}
\toprule
\textbf{Ano} & \textbf{Total} & \textbf{SUS} & \textbf{Privado} & \textbf{Estab.} & \textbf{RM/Mhab}\\
\midrule
2013 &  4.526 & 2.073 & 2.453 & 3.661 & 22,5 \\
2014 &  4.724 & 2.125 & 2.600 & 3.813 & 23,3 \\
2015 &  4.881 & 2.139 & 2.742 & 3.911 & 23,9 \\
2016 &  5.146 & 2.178 & 2.968 & 4.102 & 25,0 \\
2017 &  5.499 & 2.308 & 3.190 & 4.369 & 26,5 \\
2018 &  5.901 & 2.371 & 3.530 & 4.698 & 28,3 \\
2019 &  6.385 & 2.428 & 3.957 & 5.063 & 30,4 \\
2020 &  6.486 & 2.583 & 3.903 & 5.092 & 30,6 \\
2021 &  7.016 & 2.811 & 4.205 & 5.371 & 32,9 \\
2022 &  8.488 & 3.680 & 4.808 & 5.687 & 39,8 \\
2023 &  9.028 & 3.949 & 5.078 & 6.031 & 42,4 \\
2024 &  9.433 & 4.081 & 5.352 & 6.334 & 44,4 \\
2025 & 10.079 & 4.318 & 5.762 & 6.744 & 47,2 \\
\bottomrule
\end{tabular}
\par\noindent{\small Fonte: elaboração própria, CNES/DATASUS + IBGE.
Total = média mensal de aparelhos cadastrados no ano. Estab.\ =
estabelecimentos distintos com pelo menos um RM.}
\label{tab:evolucao-rm}
\end{table}

\subsection{Composição SUS vs. Privado}

A Figura\,\ref{fig:sus-priv} apresenta a decomposição anual da
disponibilidade de RM entre SUS e setor privado. A participação do
SUS oscila no intervalo 38\textendash 46\% ao longo da série, com
ligeira recuperação após 2020. \textbf{O setor privado detém a
maioria persistente do parque}, e a expansão pós-2020 foi
relativamente equilibrada entre os dois setores.

\begin{figure}[H]
\centering
\includegraphics[width=0.95\linewidth]{fig05-sus-vs-priv}
\caption{Distribuição SUS vs.\ Privado de aparelhos de RM, 2013\textendash 2025}
\label{fig:sus-priv}
\par\noindent{\small Fonte: elaboração própria, CNES/DATASUS.}
\end{figure}

\subsection{Distribuição geográfica}

A Figura\,\ref{fig:top10-abs} apresenta os 10 estados com maior
volume absoluto de aparelhos de RM em 2025: \textbf{São Paulo
isolado concentra 30\,\%} do parque nacional (3.054 aparelhos), e o
trio SP+RJ+MG totaliza 56\,\%. Esta concentração é consistente com
a proporção populacional desses estados, mas \textbf{NÃO é}
proporcional aos casos esperados de DP, conforme análise per capita
abaixo.

\begin{figure}[H]
\centering
\includegraphics[width=0.95\linewidth]{fig06-top10-absolute}
\caption{Top-10 UFs em número absoluto de aparelhos de RM, 2025}
\label{fig:top10-abs}
\par\noindent{\small Fonte: elaboração própria, CNES/DATASUS.}
\end{figure}

A Figura\,\ref{fig:top10-pm} apresenta a versão per capita,
revelando padrão substancialmente diferente: \textbf{Rio de Janeiro}
(82,99/Mhab), \textbf{Distrito Federal} (72,63/Mhab), \textbf{São
Paulo} (66,26/Mhab), Minas Gerais (51,86/Mhab), Rondônia
(49,06/Mhab), Paraná (47,54/Mhab) e Mato Grosso (47,26/Mhab) lideram.
O extremo oposto inclui Amazonas (15,46/Mhab), Roraima (16,24/Mhab)
e Piauí (19,20/Mhab) \textemdash{} todos abaixo da mediana OCDE.

\begin{figure}[H]
\centering
\includegraphics[width=0.95\linewidth]{fig07-top10-per-million}
\caption{Top-10 UFs em RM por 1 milhão de habitantes, 2025}
\label{fig:top10-pm}
\par\noindent{\small Fonte: elaboração própria, CNES/DATASUS + IBGE.
Linha tracejada: mediana OCDE.}
\end{figure}

A visualização integral por UF é apresentada no mapa coroplético da
Figura\,\ref{fig:choro-pm}.

\begin{figure}[H]
\centering
\includegraphics[width=0.78\linewidth]{fig08-choropleth-per-million}
\caption{Mapa coroplético — RM por milhão de habitantes por UF, 2025}
\label{fig:choro-pm}
\par\noindent{\small Fonte: elaboração própria. GeoJSON IBGE +
microdados CNES/DATASUS.}
\end{figure}

A síntese por região (Figura\,\ref{fig:region}) consolida o
gradiente Sudeste (64,82) $\rightarrow$ Centro-Oeste (49,49)
$\rightarrow$ Sul (43,20) $\rightarrow$ Nordeste (28,58)
$\rightarrow$ Norte (25,55). \textbf{O Sudeste possui mais que o
dobro de RM per capita do Norte}.

\begin{figure}[H]
\centering
\includegraphics[width=0.95\linewidth]{fig09-region-comparison}
\caption{Densidade regional de RM por milhão de habitantes, 2025}
\label{fig:region}
\par\noindent{\small Fonte: elaboração própria, CNES/DATASUS +
IBGE. Anotações internas: \% SUS por região.}
\end{figure}

\subsection{Evolução da iniquidade}

A Figura\,\ref{fig:cv} apresenta o coeficiente de variação per
capita ao longo do tempo \textemdash{} indicador de iniquidade
relativa entre UFs. \textbf{O CV permanece estável no intervalo
0,44\textendash 0,49 ao longo de toda a série}: a expansão acelerada
do parque, dobrando o número absoluto e a densidade per capita, NÃO
reduziu a iniquidade geográfica. Este achado tem implicação política
direta: políticas de expansão indiferenciada (\textit{laissez-faire})
não corrigem desigualdade estrutural; é necessária política
deliberada de focalização nos estados subdotados.

\begin{figure}[H]
\centering
\includegraphics[width=0.95\linewidth]{fig10-cv-time}
\caption{Coeficiente de variação per capita entre UFs,
2013\textendash 2025, com referência (Bolsa Família \,$\sim$\,0,45)}
\label{fig:cv}
\par\noindent{\small Fonte: elaboração própria. CV calculado sobre
as densidades RM/Mhab das 27 UFs.}
\end{figure}

A trajetória detalhada por UF \textemdash{} comparando 2013 e 2025
em densidade per capita \textemdash{} é apresentada na
Figura\,\ref{fig:growth-uf}. Todos os estados cresceram, com exceção
notável de \textbf{Roraima}, que apresentou \textit{redução}
relativa: de 22,5/Mhab em 2013 para 16,2/Mhab em 2025
(crescimento populacional superior à expansão do parque local).

\begin{figure}[H]
\centering
\includegraphics[width=0.85\linewidth]{fig11-growth-uf-2013-2025}
\caption{Crescimento por UF de RM/Mhab entre 2013 e 2025
(\textit{barbell chart})}
\label{fig:growth-uf}
\par\noindent{\small Fonte: elaboração própria. Linhas em verde =
crescimento; em vermelho = decréscimo. Apenas Roraima decresceu.}
\end{figure}

\subsection{Capacidade de neuroimagem versus carga estimada de DP}

A Figura\,\ref{fig:pd-burden} cruza, para 2025, os casos estimados
de Parkinson por UF (proxy via população $\times$ 0,33\%, baseado em
prevalência ELSI-Brazil ajustada para população 50+) com a
densidade de RM. Estados com maior carga absoluta de DP esperada
(SP, MG, RJ) tendem também a ter maior densidade de RM \textemdash{}
mas \textbf{a relação não é proporcional}: estados de carga
intermediária (BA, PE, CE) têm densidade de RM relativamente baixa
para sua população, configurando \textit{undersupply} relativo.

\begin{figure}[H]
\centering
\includegraphics[width=0.95\linewidth]{fig12-pd-burden-vs-rm}
\caption{Casos estimados de DP por UF vs.\ densidade de RM/Mhab, 2025}
\label{fig:pd-burden}
\par\noindent{\small Fonte: elaboração própria. Casos de DP estimados
como pop $\times$ 0,33\%, simplificação baseada em ELSI-Brazil
(Schlickmann \textit{et al.}, 2025); ignora variação inter-UF na
proporção 50+ ($\pm$\,2 p.p.). Linha tracejada: mediana OCDE.}
\end{figure}

% =====================================================================
\section{Discussão}\label{sec:discussao}
% =====================================================================

\subsection{Implicações dos dados epidemiológicos para estratégia
de saúde pública}

Os achados do ELSI-Brazil sobre prevalência de DP no Brasil, quando
integrados com dados atualizados de infraestrutura de saúde, revelam
lacunas críticas entre necessidade de saúde da população e capacidade
do sistema público de responder. Com projeção de 1.250.638 casos de
DP em 2060 (aumento de 133\% em relação a 2024), o Brasil enfrentará
crescimento explosivo na demanda por diagnóstico diferencial e manejo
especializado de DP\,[\ref{ref:elsi}].

\subsection{Iniquidade estrutural não resolvida pelo crescimento
agregado}

O presente trabalho documenta empiricamente um achado que tem
implicação política direta: \textbf{mesmo dobrando o parque nacional
de RM em uma década, a iniquidade per capita entre UFs não diminuiu}.
O CV permaneceu praticamente constante. Isso desmonta o argumento de
que basta investir em escala para corrigir desigualdade
\textemdash{} sem política focalizada (incentivo fiscal regional,
financiamento priorizado pelo Ministério da Saúde, parceria
público-privada com cláusulas de localização), o crescimento
agregado se reproduz na distribuição existente.

\subsection{Potencial de neuroimagem em aprimoramento diagnóstico:
perspectivas para integração no SUS}

Embora neuroimagem seja frequentemente vista como luxo diagnóstico
inadequado em contextos de recursos limitados, argumento pode ser
construído para sua integração seletiva em protocolo diagnóstico de
DP em Brasil. Especificamente, utilização de RM estrutural em casos
de apresentação clínica ambígua (tremor predominante versus rigidez,
ou quando diagnóstico diferencial inclui parkinsonismos atípicos)
ofereceria valor diagnóstico agregado que justificaria alocação de
recursos limitados\,[\ref{ref:saba}]. Modelo implementado em
contextos similares de países de renda média (por exemplo, México,
Argentina) demonstrou que centralização de protocolos de RM em
centros de referência regional, combinada com educação de clínicos
locais em reconhecimento de características diagnósticas de DP,
melhora acurácia diagnóstica em 25\textendash 30\% comparado a
cenário onde apenas diagnóstico clínico está
disponível\,[\ref{ref:datasus}].

Para implementação no SUS especificamente, recomendação seria:
\textbf{(1)} desenvolver protocolo padronizado de RM para DP e
parkinsonismos atípicos que pudesse ser implementado em equipamentos
clínicos já existentes sem requerimento de tecnologia adicional;
\textbf{(2)} capacitar radiologistas em interpretação especializada
de imagens neurológicas focadas em distúrbios de movimento;
\textbf{(3)} criar centros de referência regional em cada região
geopolítica brasileira onde casos complexos de apresentação ambígua
poderiam ser encaminhados para confirmação diagnóstica com
neuroimagem e avaliação neurológica especializada.

Custo estimado para desenvolvimento de protocolo padronizado e
capacitação inicial de radiologistas seria de aproximadamente
\rs 500.000 a \rs 1.000.000 em investimento único, com custos
operacionais recorrentes de exames adicionais de RM sendo absorvíveis
dentro de orçamentos existentes de equipamentos já
instalados\,[\ref{ref:datasus}].

\subsection{Limitações da inferência}

\textbf{Primeira}, os resultados são predominantemente descritivos:
nenhum desenho identifica relação causal isolada entre disponibilidade
de RM e desfechos clínicos de DP. \textbf{Segunda}, o ELSI-Brazil
captura DP autorrelatada e amostra preferentemente indivíduos em
estágios avançados (29\% precisam de suporte para caminhar) \textemdash{}
subestima prevalência inicial. \textbf{Terceira}, contagem CNES não
distingue equipamento operacional vs.\ cadastrado; cruzamento futuro
com produção SIH-RD/SIA-RD daria visão mais precisa.

% =====================================================================
\section{Considerações Finais}\label{sec:conclusao}
% =====================================================================

\subsection{O que descobrimos}

\textbf{(i) O Brasil dispõe, em 2025, de \,$\sim$\,10.079 aparelhos
de Ressonância Magnética em 6.744 estabelecimentos}, com
participação do SUS em torno de 43\%. A densidade nacional alcançou
47,2/Mhab, acima da mediana OCDE.

\textbf{(ii) A iniquidade geográfica permanece estrutural.} CV
estável em 0,44\textendash 0,49 ao longo de 12 anos. RJ-DF-SP no
topo (66\textendash 83/Mhab) vs.\ AM-RR-PI no fundo
(15\textendash 19/Mhab). Crescimento agregado não resolve
desigualdade.

\textbf{(iii) Carga epidemiológica crescente.} Projeção de 1,25
milhão de casos de DP no Brasil até 2060, com gradiente
distribuído por todas as regiões. Necessidade de diagnóstico
especializado crescerá no mínimo proporcionalmente.

\textbf{(iv) Neuroimagem é complementar, não substitutiva, ao
diagnóstico clínico.} Mas em casos ambíguos, RM estrutural
adequadamente protocolizada pode aumentar acurácia diagnóstica em
25\textendash 30\%, segundo evidência de países de renda média.

\textbf{(v) A solução é organizacional, não tecnológica.} O parque
de equipamentos já existe; o gap está em protocolos padronizados,
capacitação de radiologistas em distúrbios de movimento, e
infraestrutura de teleneurologia para regiões periféricas.

\subsection{Agenda de pesquisa}

Quatro frentes destacam-se: \textbf{(a)} cruzamento de CNES
(equipamentos) com SIH-AIH (produção) para mensurar utilização
\textit{efetiva} dos aparelhos por UF; \textbf{(b)} cruzamento com
microdados de SIM-DO para óbitos com causa associada à DP, validando
a estimativa de carga; \textbf{(c)} estudo de validação de
protocolos padronizados de RM em contextos brasileiros, conforme
recomendado por Rolim (2026); \textbf{(d)} extensão da janela
temporal para 2005\textendash 2012 (já materializada na bronze) para
capturar período de baixa cobertura inicial do CNES e identificar
inflexões.

\subsection{Contribuição metodológica}

Para além dos achados substantivos, o presente artigo demonstra um
modelo de \textbf{coautoria entre pesquisa clínica-epidemiológica
(Rolim) e engenharia de dados (Chalhoub)}: a primeira fornece o
arcabouço teórico, conceitual e propositivo; a segunda atualiza,
formaliza e torna reproduzíveis as quantificações que sustentam o
argumento. Acreditamos que esse modelo \textemdash{} cuja viabilidade
é demonstrada pela presente integração entre o manuscrito original
de Rolim e o pipeline aberto \textit{Mirante dos Dados} \textemdash{}
pode ser replicado para outras agendas de saúde
pública nas quais existe um manuscrito clínico-epidemiológico com
recomendações que se beneficiariam de quantificação atualizada,
auditável, e independente de revisão.

% =====================================================================
\section*{Referências}\label{sec:ref}
\addcontentsline{toc}{section}{Referências}
% =====================================================================
\begin{singlespace}\small

\noindent\hangindent=1.25cm\hangafter=1\label{ref:elsi}
[1] SCHLICKMANN, T. H.; TESSARI, M. S.; BORELLI, W. V. \textit{et al.}
Prevalence, distribution and future projections of Parkinson disease
in Brazil: insights from the ELSI-Brazil cohort study. \textbf{Lancet
Reg Health Am.}, v. 44, n. 101046, 2025. DOI:
10.1016/j.lana.2025.101046.\par\smallskip

\noindent\hangindent=1.25cm\hangafter=1\label{ref:su}
[2] SU, D. \textit{et al.} Projections for prevalence of Parkinson's
disease and its driving factors in 195 countries and territories to
2050: modelling study of Global Burden of Disease Study 2021.
\textbf{BMJ}, 2025. DOI: 10.1136/bmj-2024-080952.\par\smallskip

\noindent\hangindent=1.25cm\hangafter=1\label{ref:saba}
[3] SABA, R. A.; MAIA, D. P.; CARDOSO, F. E. C. \textit{et al.}
Guidelines for Parkinson's disease treatment: consensus from the
Movement Disorders Scientific Department of the Brazilian Academy
of Neurology — motor symptoms. \textbf{Arq Neuropsiquiatr.}, v. 80,
n. 3, p. 316\textendash 329, 2022. DOI: 10.1590/0004-282X-ANP-2021-0219.
\par\smallskip

\noindent\hangindent=1.25cm\hangafter=1\label{ref:bovolenta}
[4] BOVOLENTA, T. M.; SCHUMACHER-SCHUH, A. F.; SANTOS-LOBATO, B. L. D.
\textit{et al.} Average annual cost of Parkinson's disease in a
Brazilian multi-ethnic population. \textbf{Parkinsonism Relat
Disord.}, v. 117, n. 105897, 2023. DOI:
10.1016/j.parkreldis.2023.105897.\par\smallskip

\noindent\hangindent=1.25cm\hangafter=1\label{ref:datasus}
[5] DEPARTAMENTO DE INFORMÁTICA DO SISTEMA ÚNICO DE SAÚDE (DATASUS).
\textbf{Recursos de Diagnóstico por Imagem no SUS}. Ministério da
Saúde do Brasil, 2024. Disponível em:
\url{https://datasus.saude.gov.br/}.\par\smallskip

\noindent\hangindent=1.25cm\hangafter=1
[6] INSTITUTO BRASILEIRO DE GEOGRAFIA E ESTATÍSTICA (IBGE).
\textbf{Projeções da População do Brasil: 2010\textendash 2060}.
2013.\par\smallskip

\noindent\hangindent=1.25cm\hangafter=1
[7] LLIBRE-GUERRA, J. J.; PRINA, M.; SOSA, A. L. \textit{et al.}
Prevalence of parkinsonism and Parkinson disease in urban and rural
populations from Latin America: a community based study.
\textbf{Lancet Reg Health Am.}, v. 7, n. 100136, 2022. DOI:
10.1016/j.lana.2021.100136.\par\smallskip

\noindent\hangindent=1.25cm\hangafter=1
[8] MARRAS, C.; BECK, J. C.; BOWER, J. H. \textit{et al.}
Prevalence of Parkinson's disease across North America. \textbf{NPJ
Parkinsons Dis.}, v. 4, n. 1, p. 21, 2018. DOI:
10.1038/s41531-018-0058-0.\par\smallskip

\noindent\hangindent=1.25cm\hangafter=1
[9] BEN-SHLOMO, Y.; DARWEESH, S.; LLIBRE-GUERRA, J. \textit{et al.}
The epidemiology of Parkinson's disease. \textbf{Lancet}, v. 403,
n. 10423, p. 283\textendash 292, 2024. DOI: 10.1016/S0140-6736(23)01419-8.
\par\smallskip

\noindent\hangindent=1.25cm\hangafter=1
[10] TOLOSA, E.; GARRIDO, A.; SCHOLZ, S. W.; POEWE, W. Challenges
in the diagnosis of Parkinson's disease. \textbf{Lancet Neurol.},
v. 20, n. 5, p. 385\textendash 397, 2021. DOI:
10.1016/S1474-4422(21)00030-2.\par\smallskip

\noindent\hangindent=1.25cm\hangafter=1
[11] FASANO, A.; SCHNEIDER, S. A.; ABRAMYCHEVA, N. Y. \textit{et al.}
Predictive factors for deep brain stimulation outcome in
tremor-dominant Parkinson disease. \textbf{Neurology}, v. 97, n. 5,
p. e520\textendash e531, 2021. DOI: 10.1212/WNL.0000000000012358.
\par\smallskip

\noindent\hangindent=1.25cm\hangafter=1
[12] PRINGSHEIM, T.; JETTE, N.; FROLKIS, A.; STEEVES, T. D. L. The
prevalence of Parkinson's disease: a systematic review and
meta-analysis. \textbf{Mov Disord.}, v. 29, n. 13,
p. 1583\textendash 1590, 2014. DOI: 10.1002/mds.25945.\par\smallskip

\noindent\hangindent=1.25cm\hangafter=1
[13] DORSEY, E. R.; ELBAZ, A.; NICHOLS, E. \textit{et al.} Global,
regional, and national burden of Parkinson's disease,
1990\textendash 2016: a systematic analysis for the Global Burden
of Disease Study 2016. \textbf{Lancet Neurol.}, v. 17, n. 11,
p. 939\textendash 953, 2018. DOI: 10.1016/S1474-4422(18)30295-3.
\par\smallskip

\noindent\hangindent=1.25cm\hangafter=1
[14] SCHIESS, N.; CATALDI, R.; OKUN, M. S. \textit{et al.} Six
action steps to address global disparities in Parkinson disease: a
World Health Organization priority. \textbf{JAMA Neurol.}, v. 79,
n. 9, p. 929\textendash 936, 2022. DOI: 10.1001/jamaneurol.2022.1783.
\par\smallskip

\noindent\hangindent=1.25cm\hangafter=1
[15] POSTUMA, R. B.; BERG, D.; STERN, M. \textit{et al.} MDS
clinical diagnostic criteria for Parkinson's disease.
\textbf{Mov Disord.}, v. 30, n. 12, p. 1591\textendash 1601, 2015.
DOI: 10.1002/mds.26424.\par\smallskip

\noindent\hangindent=1.25cm\hangafter=1
[16] ARMBRUST, M.; GHODSI, A.; XIN, R.; ZAHARIA, M. Lakehouse: a
new generation of open platforms that unify data warehousing and
advanced analytics. In: \textbf{11th CIDR}, 2021.\par\smallskip

\noindent\hangindent=1.25cm\hangafter=1
[17] WILKINSON, M. D. \textit{et al.} The FAIR Guiding Principles
for scientific data management and stewardship. \textbf{Scientific
Data}, v. 3, 2016. DOI: 10.1038/sdata.2016.18.\par\smallskip

\noindent\hangindent=1.25cm\hangafter=1
[18] ROLIM, A. M. Análise estatística da prevalência da Doença de
Parkinson no Brasil: desafios, tecnologias de neuroimagem e
perspectivas de políticas públicas. \textbf{Manuscrito v1}, abr.
2026. (Texto-base deste Working Paper, integrado nas seções 1\textendash 2.)
\par\smallskip
\end{singlespace}

\vspace{1cm}

\noindent\textbf{Citar como:}\\
ROLIM, A. M.; CHALHOUB, L. Análise estatística da prevalência da
Doença de Parkinson no Brasil: desafios, tecnologias de neuroimagem
e perspectivas de políticas públicas.
\textit{Mirante dos Dados \textemdash{} Working Paper n. 4}, v. 1.0,
abr. 2026. Disponível em:
\url{https://leonardochalhoub.github.io/mirante-dos-dados-br/}.

\vspace{0.5cm}

\noindent\textbf{Reconhecimento:} a análise epidemiológica e a
revisão de literatura que sustentam este artigo foram originalmente
escritas por Alexandre Maciel Rolim como manuscrito independente
(v1, abr/2026). A presente versão integra esse manuscrito com a
camada de engenharia de dados do \textit{Mirante dos Dados},
contribuindo análise atualizada de infraestrutura de RM via
microdados CNES 2013\textendash 2025 e formalização do pipeline
reproduzível.

\vspace{0.5cm}

\noindent\textbf{Licença:} dados consolidados (camada Gold) e código
do pipeline distribuídos sob licença MIT. Texto deste artigo
distribuído sob licença Creative Commons Atribuição 4.0 Internacional
(CC BY 4.0).

\end{document}
